\section{Introduction}
\label{sec:introduction}

A used Quadro RTX 6000 sells for roughly a tenth of its launch price, and
two of them bridged by NVLink put 48~GB of GPU memory and a measured
1.2~TB/s of aggregate DRAM bandwidth on a desk. Whether that hardware can
carry a modern inference workload is not answerable from marketing
geometry: it depends on instruction latencies, cache behaviour, and
interconnect costs that NVIDIA does not publish and that existing
microbenchmark studies cover only for neighbouring chips. Jia et
al.\ dissected the TU104-based Tesla T4 \citep{jia2019turing}; the larger
TU102, with its different L2, DRAM system, and NVLink, has no equivalent
public characterisation.

This paper measures one. In the style of Agner Fog's x86 instruction
tables \citep{fog2025instruction}, it builds a complete operation table
for the TU102 at compute capability 7.5 (instruction latency and
throughput with measured issue-pipe bindings, the full memory hierarchy
from the register file to DRAM, and the NVLink, PCIe, and NCCL
interconnect costs) from microbenchmarks whose every timed loop is
verified at the SASS level and whose every published row carries its
provenance.

The research question is operational rather than encyclopaedic:
\emph{can a table of independently measured primitive costs predict the
behaviour of real composite kernels well enough to drive engineering
decisions?} The answer, developed through three pre-registered hypotheses
(\cref{sec:hypotheses}), is sharply split. Composed predictions are
excellent where a single resource binds, within 3.5\% on
single-resource kernels and within 10\% on a latency-dominated
interconnect primitive, and fail systematically where work couples across issue
pipes or where a primitive row turns out to be a floor rather than a
typical cost. Both failure modes are measured, diagnosed, and published.

The contributions are: (i) the table itself, 264 rows with per-row
provenance, priors, and deviation flags, regenerable from a fresh clone;
(ii) the measurement methodology, including a SASS-gated benchmark
discipline that caught every compiler-induced measurement defect before a
number shipped, and a clock-domain contract for a GPU whose memory clock
cannot be locked; (iii) findings absent from public documentation, among
them the fma-pipe binding of \texttt{IDP.4A}, the full-rate
f32-accumulate of Quadro-positioned TU102 tensor cores, the L1-retention
semantics of \texttt{.cs} loads, and the request-path saturation of the
L2; and (iv) a worked demonstration that pre-registration (gates,
metrics, and decision rules committed to a public repository before data
collection) transfers from empirical software engineering to
microbenchmarking, where the temptation to tune until numbers match
priors is the central threat to validity.
